\documentstyle[%


righttag%


,11pt%


,amssymb%


,russian%               remove in Eng.


,izvru%                 Set izven% in Eng.


% ,izvru-ks             for brief communications


]{amsart}





%





\newcommand{\Mod}{\operatorname{mod}}


\newcommand{\tts}[1]{}





\theoremstyle{plain}


\theorembodyfont{\it}


\newtheorem{props}{\hskip\parindent Предложение}[section]





\theoremstyle{definition}


\theorembodyfont{\rm}


\newtheorem{cors}{\hskip\parindent Следствие}[section]


\newtheorem{notenonum}{\hskip\parindent Замечание}


\renewcommand{\thenotenonum}{}





%\hoffset=-15mm%                  For Eng.


\setcounter{page}{79}


\god{2001}


\nomer{2~(465)} %                        Remove in Eng.


%\nomer{Vol.\,45, No.\,2}%               Set in Eng.


%\str{\pageref{begin}--\pageref{end}}%  Set in Eng.


\udk{517.584}





\author{\it С.П.\,Хэкало}


% NB:   ^^ remove in Eng.


\title{Функция Бесселя на конечном поле}


%\thanks{Работа выполнена при поддержке Российского фонда


%фундаментальных исследований, грант \No\, ...}





\date{}


%\pagestyle{myheadings}%         Set in Eng.


\pagestyle{plain} %              Remove in Eng.


\begin{document}


%\thispagestyle{plain}%          Set in Eng.


\maketitle


\label{begin}





Аналоги бесселевых функций с аргументом из конечного поля 


возникают как в теории


специальных функций [1], [2], так и в теории представлений групп [3]--[6].


В данной работе изучается аналог функции Бесселя, через который


выражаются матричные элементы неприводимых представлений группы


движений плоскости над конечным полем. Метод, развитый в [7],


позволяет получить формулу Хансена,


теоремы сложения и умножения.





\section{Определение функции Бесселя}





Пусть $\Bbb F$ --- конечное поле из $q=p^n$ элементов


($p$ --- характеристика поля), $\Bbb F^*$ --- мультипликативная группа поля,


$\Bbb F(\tau)=\{z\}$ --- квадратичное расширение


с сопряжением $\ov{z}$ и нормой


$N(z)=z\ov{z}$. Через $\eta$ обозначим образующую циклической группы


$\Bbb F^*(\tau)$. Для каждого $u\in{\Bbb F^*}$


окружность $C_u=\{z\in{\Bbb F(\tau)}\mid


N(z)=u\}$ состоит из $q+1$ точек. Элемент


$\xi=\eta^{q-1}$ порождает циклическую


группу $C\equiv{C_1}$ единиц поля $\Bbb F(\tau)$


[8]. Пусть $k(t)$ --- показатель степени в равенстве


$t=\xi^{k(t)}$, $t\in{C}$. Как известно (напр., [8]), характеры


группы $C$ имеют вид


$$


\pi_s(t)=\exp \bigg(\frac{2\pi i}{q+1}sk(t) \bigg),\quad s=0,1,\dotsc,q,


$$


где $s$ --- номер характера.


Введем билинейную форму $\langle z_1,z_2\rangle=z_1z_2+\ov{z}_1\ov{z}_2$


и зафиксируем нетривиальный аддитивный характер


$\chi$ поля $\Bbb F$. Бесселеву функцию


$J_s(z)$ аргумента $z\in{\Bbb F(\tau)}$


и индекса $s=0,1,\dotsc,q$ со значениями в поле


$\Bbb C$ комплексных чисел определим равенством


$$


J_s(z):=\frac{1}{q+1} \sum_{t,N(t)=1}{\chi(\langle z,t\rangle)\pi_s(t)}.


$$





Близкая к $J_s(z)$ функция рассматривалась в [4]


\begin{equation}


j_r(u)=\frac{1}{q}\sum_{t, N(t)=u}{\chi(t+\ov{t})\varphi_r(t)},


\end{equation}


где $r=0,1,\dotsc,q^2-2$, $\varphi_r(t)$ ---


мультипликативный характер поля $\Bbb F^*(\tau)$.


Пусть ${[r]}\equiv{r(\Mod(q+1))}$ и ${0}\le{[r]}\le{q}$,


тогда связь между указанными функциями такова:


\begin{equation}


j_r(z\ov{z})=\frac{q+1}{q}\varphi_r(z)J_{[r]}(z).


\end{equation}





\section{Свойства функции $J_s(z)$}





Непосредственно из определения следует, что


для $z\in{\Bbb F(\tau)}$ и $ s=0,1,\dotsc,q$


\begin{gather*}


\sum_{z\in{\Bbb F(\tau)}}{J_s(z)}=0;\quad


J_s(-z)=(-1)^sJ_s(z), \quad p\neq2,\\


\ov{J}_s(z)=J_s(-\ov{z}); \quad


J_{[-s]}(z)=J_s(\ov{z}).


\end{gather*} 





Связь (2) и тождества для функции (1) из [4] порождают следующие тождества


для функции $J_s(z)$.





\begin{props}


Для $[r]=1,2,\dotsc,q$


\begin{align*}


&1. \quad \sum_{z\in{\Bbb F^*(\tau)}}{\varphi_r(z\ov{z}^{-1})J_{[r]}^2(z)}=


\frac{q^2}{q+1}


\begin{cases} (-1)^{[r]}, & p\neq2;\\ 1, & p=2;\end{cases} \\


%


&2.\quad \sum_{z\in{\Bbb F^*(\tau)}}{\varphi_r(z\ov{z}^{-1})J_{[r]}(tz)


J_{[r]}(z)}=0,\quad N(t)\neq1;\\


%


&3.\quad \sum_{z\in{\Bbb F^*(\tau)}}{\varphi_r(z\ov{z}^{-1}){\chi(z\ov{z})}


J_{[r]}(xz)J_{[r]}(yz)}


=-q\chi(-x\ov{x}-y\ov{y})J_{[r]}(xy)


\begin{cases} (-1)^{[r]}, & p\neq2;\\ 1, & p=2.\end{cases}


\end{align*}


\end{props}





Дополнительные свойства функции $J_s(z)$ будут получены методом 


теории представлений групп [7].





\section{Матричные элементы}





Движением плоскости $\Bbb F(\tau)$ будем


называть биекции типа $w=hz+a$, где $a\in{\Bbb F(\tau)}$ задает перенос, а


$h=\alpha+\tau\beta$ при $N(h)=1$ задает вращение.





Приведем необходимые сведения о группе $G$ всех таких движений плоскости.


Группа $G$ является полупрямым произведением


$\Bbb F(\tau)\rtimes{C}$ коммутативного


нормального делителя $\Bbb F(\tau)$


переносов и циклической группы $C$ вращений. 


Она отождествляется с группой матриц


$$


g(h,a)=\begin{pmatrix} h & a \\ 0 & 1 \end{pmatrix},\quad


h \in C;\quad a\in {\Bbb F(\tau)}.


$$





Таким образом, порядок группы $G$ равен $|G|=q^2(q+1)$.





Подгруппа $C$ группы $G$ действует на окружности $C_u$, $u\in{\Bbb F^*},$


транзитивно и на $\Bbb F(\tau)$ самосопряженно относительно скалярного


произведения $\langle\cdot,\cdot\rangle.$


Неприводимые представления группы $G$ строятся в соответствии


с общей теорией Макки (напр., [9], с.\,63).





Обозначим через $\varPhi=\{f\mid f: C\rightarrow{\Bbb C}\}$ пространство


комплекснозначных функций на окружности $C$.


Неприводимое представление $T(g): \varPhi\rightarrow\varPhi$ действует по


формуле


$$


T(g(h,a))f(z)=\chi(\langle z,a\rangle)f(hz).


$$





Указанное представление интересно, поскольку его матричные элементы


в базисе из характеров на $C$ имеют  вид


\begin{equation}


t_{nm}(g)=\pi_m(h)J_{[m-n]}(a),


\end{equation}


где


$0\le n,m\le q$, $h\in C$; $\pi$ --- мультипликативный характер на $C$.





Равенство (3) является основным для 


вывода свойств бесселевой функции $J_s(z)$.





\section{Дополнительные свойства}





В этом пункте приведем дополнительные свойства


функции $J_s(z)$ (теоремы сложения, умножения,


формулу Хансена и др.), вытекающие


из теории представлений. Обозначим дельта-функцию через $\delta$.





\begin{props}


$J_s(0)=\delta(s)$.


\end{props}





\begin{props}


$\displaystyle\sum\limits_{z\in {\Bbb F(\tau)}}{J_0^2(z)=\frac{q^2}{q+1}}.$


\end{props}





\begin{props}[{\series{m}\shape{n}\selectfont теорема сложения}]


$$


J_s(z_1+hz_2)=\sum_{j=0}^q{\pi_{[j-s]}(h)J_j(z_1)J_{[s-j]}(z_2)},


$$


где $z_1,z_2\in{\Bbb F(\tau)}$ и $h\in C$.


\end{props}





\begin{cors}


\begin{gather*}


J_s(z_1+z_2)=\sum_{j=0}^q{J_j(z_1)J_{[s-j]}(z_2)};\quad


J_s(hz)={\ov{\pi}_s(h)J_s(z)};\\


J_0(z_1+hz_2)=\sum_{j=0}^q{{\pi_j}(h)J_j(z_1)J_j(\ov{z}_2)};\quad


J_0(z_1+z_2)=\sum_{j=0}^q{J_j(z_1)J_j(\ov{z}_2)},


\end{gather*}


где $z_1,z_2\in{\Bbb F(\tau)}$ и $h\in C$.


\end{cors}





\begin{props}[{\series{m}\shape{n}\selectfont формула Хансена}]


\begin{alignat*}{2}


\sum_{j=0}^q{(-1)^jJ_j(z)J_{[s-j]}(z)}&=\delta(s),&\quad p&\neq2;\\


\sum_{j=0}^q{J_j(z)J_{[s-j]}(z)}&=\delta(s), &\quad p&=2.


\end{alignat*}


\end{props}





\begin{props}[{\series{m}\shape{n}\selectfont теорема умножения}]


$$


\frac{1}{q+1}\sum_{h, N(h)=1}{J_s(z_1+hz_2)}=J_s(z_1)J_0(z_2).


$$


\end{props}





\begin{cors}


$$


\frac{1}{q+1}\sum_{h, N(h)=1}{J_s(hz)}=J_0(z)\delta(s).


$$


\end{cors}





\begin{notenonum}


Равенства п.\,2, а также предложения 4.1--4.5 могут быть 


использованы для нахождения свойств функции $j_r(u)$,


например, аналогом предложения


4.3 является равенство


$$


\ov\varphi_r(z)j_r(z\ov{z})=\frac{q}{q+1}\sum_{n=0}^q{\varphi_{n-r}(hz_2)}


{\ov\varphi_n(z_1)


j_n(z_1\ov{z}_1) j_{r-n}(z_2\ov{z}_2)},


$$


где $z=z_1+hz_2$ и $\varphi$ ---


мультипликативный характер на $\Bbb F^*(\tau)$.


\end{notenonum}





\section{Ограничение функции Бесселя на поле $\Bbb F$}





На основе свойств функции $J_s(z)$, 


определенной на расширении $\Bbb F(\tau)$


(см. пп.\,2,~4), дадим список соответствующих равенств для функции Бесселя


$J_s(x)$, определенной на исходном поле~${\Bbb F}.$ 





\begin{props}


\begin{gather*}


\sum_{x\in{\Bbb F}}{J_s(x)}=0;\quad


J_s(-x)=(-1)^sJ_s(x), \quad p\neq2;\\


%


\ov{J}_s(x)=J_s(-x);\quad


J_{[-s]}(x)=J_s(x);\quad J_s(0)=\delta(s). \\


%


J_s(x_1+x_2)=\sum_{j=0}^q{J_j(x_1)J_{[s-j]}(x_2)};\\


%


\sum_{j=0}^q{(-1)^jJ_j(x)J_{[s-j]}(x)}=\delta(s),\quad p\neq2;\\


%


\sum_{j=0}^q{J_j(x)J_{[s-j]}(x)}=\delta(s),\quad p=2.


\end{gather*}


\end{props}





\begin{props}


$$


\sum_{x\in{\Bbb F}}{J_0^2(x)}=


\begin{cases} 2q^2(q+1)^{-2}, & p\neq2; \\


q(2q+1)(q+1)^{-2}, & p=2.


\end{cases}


$$


\end{props}





Автор благодарен Е.Е.\,Петрову и В.П.\,Лексину за полезные советы


и критические замечания.





\begin{thebibliography}{99}





\bibitem{1}


Evans J. {\em Hermite character sums} // Pacific J. of


Math. -- 1986. -- V.\,122. --  \No\,2. -- P.\,375--390.


\bibitem{2}


Greene J. {\em Hypergeometric functions over finite fields}


// Trans. Amer. Math. Soc. -- 1987. -- V.\,301. -- \No\,1.


-- P.\,77--101.


\bibitem{3}


Гельфанд С.И. {\em Представление полной линейной группы над конечным полем}


// Матем. сб. -- 1970. -- Т.\,83. -- \No\,9. -- С.\,15--41.


\bibitem{4}


Piatetski-Shapiro J. {\em Complex representation of $GL(2,k)$ for finite


fields $k$} // Contemp. Math. -- 1987. -- V.\,16. -- P.\,1--71.


\bibitem{5}


Vigneras M.-F. {\em Representations modulaires


de $GL(2,F)$ en caracteristique


$l$, $F$ corps fini de caracteristique $p\neq l$} // C.R. Acad. Sci., Paris.


-- 1998. -- T.\,306. -- Ser.~I. -- P.\,451--454.


\bibitem{6}


Soto-Andrade J. {\em Geometrical Gel'fand models, tensor quotients and


Weil representation} // Proc. Symp. Pure Math. -- 1987.


-- V.\,48. -- P.\,305--316.


\bibitem{7}


Виленкин Н.Я. {\em Специальные функции и теория представлений групп}. --


М.: Наука, 1965. -- 588~с.


\bibitem{8}


Лидл Р., Нидеррайтер Г. {\em Конечные поля}. T.\,1. -- М.: Наука, 1988.


-- 428~c.


\bibitem{9}


Серр Ж.-П. {\em Линейные представления конечных групп}.


-- М.: Мир, 1970. -- 132~с.








\end{thebibliography}





\vskip 2cm





\post{\it Коломенский педагогический институт}{\it Поступила}


\post{}{17.08.1998}





\label{end}


\end{document}





